\section{Related Work}

\subsection{Skin Lesion Synthesis}
Generative models have been applied to dermoscopy synthesis since the advent of Generative Adversarial Networks (GANs). Early works used DCGAN and ProGAN architectures to oversample minority classes in ISIC datasets, demonstrating modest improvements in classifier accuracy~\cite{TexasGAN2023}. StyleGAN2 and its adaptive discriminator augmentation variant (StyleGAN2-ADA) achieved more stable training on small medical datasets, producing high-resolution lesion images with realistic textures~\cite{SkinGenBench2025}. Conditional extensions, such as class-conditional GANs and RedGAN~\cite{qasim2021redganclassimbalance}, targeted class imbalance directly by conditioning generation on lesion type labels. More recently, diffusion models have replaced GANs as the dominant paradigm. DermDiff~\cite{DermDiff2025} adapted Stable Diffusion with text conditioning on skin tone and lesion type, while DiT-based approaches~\cite{DermDiT2025} demonstrated low FID scores on curated ISIC subsets. Concurrent work by She et al.~\cite{SheJun_DiDGen_MICCAI2025} and Xu et al.~\cite{xu2025skindualgenpromptdrivendiffusionsimultaneous} explored simultaneous image-mask generation to bypass the need for pre-existing segmentation annotations. Mask2Derm differs from these works by conditioning generation on explicit binary masks and targeting a harmonized multi-source corpus rather than a single curated dataset.

\subsection{Conditional Image Generation}
Conditioning generative models on spatial signals has a rich history. Pix2pix~\cite{isola2017image} established paired image-to-image translation via conditional GANs, while SPADE~\cite{park2019semantic} introduced spatially-adaptive normalization for semantic layout synthesis. ControlNet~\cite{zhang2023adding} extended this idea to large-scale diffusion models by injecting spatial conditioning through locked-copy encoder blocks and zero-initialized convolutions, preserving the pre-trained backbone's generative capacity while learning new conditioning pathways. InstructPix2Pix~\cite{brooks2023instructpix2pix} and IP-Adapter demonstrated that additional modalities can be integrated into diffusion models without catastrophic forgetting. We adopt ControlNet as our conditioning mechanism because it offers a favorable trade-off between geometric control fidelity and training stability on domain-shifted medical data.

\subsection{Data Augmentation for Medical Imaging}
Synthetic data augmentation has shown consistent benefits across medical imaging modalities. In radiology, GAN-generated chest X-rays and CT slices have been used to augment small annotated sets for pneumonia and lung nodule detection~\cite{TODO}. In pathology, whole-slide image synthesis has improved rare-subtype classification. For dermoscopy specifically, Du et al.~\cite{du2023boostingdermatoscopiclesionsegmentation} demonstrated that diffusion-generated images can boost lesion segmentation performance when mixed with real data. The TSTR protocol, formalized by Esteban et al.~\cite{esteban2017real}, provides a rigorous framework for evaluating whether synthetic data carries sufficient clinical signal to substitute for real annotations; we adopt this protocol as our primary utility metric.
